\documentclass{article}
\usepackage{xeCJK, amsmath, amssymb, geometry}
\setCJKmainfont{標楷體}
\geometry{a4paper, margin=2.5cm}

\title{數統導論作業四}
\author{111020003 林群佑}
\date{}

\begin{document}

\maketitle
\section*{3.1.12}
\subsection*{(a)}
\begin{align*}
M_{X_2, \ldots, X_{k-1}}(t_2, \ldots, t_{k-1}) 
&= M_{X_1, \ldots, X_{k-1}}(0, t_2, \ldots, t_{k-1}) \\
&= (p_1 e^0 + p_2 e^{t_2} + \cdots + p_{k-1} e^{t_{k-1}} + p_k)^n \\
&= (p_1 + p_k + p_2 e^{t_2} + \cdots + p_{k-1} e^{t_{k-1}})^n
\end{align*}

\subsection*{(b)}

By mgf of (a) we can know that
pmf of $X_2, \ldots, X_{k-1}$ is
\[ 
P(X_2=x_2, \ldots, X_{k-1}=x_{k-1}) = 
\frac{n!}{x_2! \cdots x_{k-1}! (n - S)!} p_2^{x_2} \cdots p_{k-1}^{x_{k-1}} (p_1+p_k)^{n-s}
\]
where $S = \sum_{i=2}^{k-1} x_i \leq n$。

\subsection*{(c)}
By definition
\begin{align*}
P(X_1=x_1 | X_2=x_2, \ldots, X_{k-1}=x_{k-1})
&= \frac{P(X_1=x_1, X_2=x_2, \ldots, X_{k-1}=x_{k-1})}{P(X_2=x_2, \ldots, X_{k-1}=x_{k-1})}\\
&=\frac{\frac{n!}{x_1! x_2! \cdots x_{k-1}! x_k!} p_1^{x_1} p_2^{x_2} \cdots p_{k-1}^{x_{k-1}} p_k^{x_k}}
{\frac{n!}{x_2! \cdots x_{k-1}! (x_1+x_k)!} p_2^{x_2} \cdots p_{k-1}^{x_{k-1}} (p_1+p_k)^{(x_1+x_k)}}\\
&=\frac{(x_1 + x_k)!}{x_1! x_k!}\frac{p_1^{x_1}p_k^{x_k}}{(p_1+p_k)^{(x_1+x_k)}}\\
&=\frac{(n-S)!}{x_1! (n-S-x_1)!}\left(\frac{p_1}{p_1+p_k}\right)^{x_1}
\left(\frac{p_k}{p_1+p_k}\right)^{n-S-x_1}
\end{align*}
where \(S = \sum_{i=2}^{k-1} x_i\)
\subsection*{(d)}
By (c) we can know that \(X_1 | X_2=x_2, \ldots, X_{k-1}=x_{k-1}
\sim \text{Binomial}(n-\sum_{i=2}^{k-1}x_i \, , \frac{p_1}{p_1+p_k})\)

\noindent
Therefore,
\[
E(X_1 | X_2 = x_2, \ldots, X_{k-1}) = 
\left( n - \sum_{i=2}^{k-1} x_i \right) \left( \frac{p_1}{p_1+p_k} \right)
\]
\section*{3.1.17}
First we compute mgf of \(\text{Geometric}(p)\) only consider fail times before success.
\begin{align*}
\mathbb{E}[e^{tx}]&=\sum_{x=0}^{\infty} e^{tx}\left(1-p\right)^x p\\
&=p\sum_{x=0}^{\infty}\left[e^t\left(1-p\right)\right]^x\\
&=\frac{p}{1-\left(1-p\right)e^t}
\end{align*}
\noindent
Because it need converge that \(t < -\ln(1-p) \)

\noindent
Then in fact that we can seen \(\text{Negative Binomial}(r, p)\) as \(r\) times independent \(\text{Geometric}(p)\).

\noindent
Therefore,
\begin{align*}
M_{NB(r,p)}(t) &= \prod_{i=1}^{r}M_{Geo(p)}(t)\\
&=\left[\frac{p}{1-\left(1-p\right)e^t}\right]^r\\
&=p^r\left[1-\left(1-p\right)e^t\right]^{-r}
\end{align*}
\noindent
It also need satisfy that \(t < -\ln(1-p) \)

\noindent
Expection of Negative Binomial is:
\begin{align*}
\mathbb{E}[X] &= \frac{d}{dt}M(0)\\
&= \frac{d}{dt} \, p^r\left[1-\left(1-p\right)e^t\right]^{-r} \\
&= p^r \cdot (-r) \cdot \left[1-\left(1-p\right)e^t\right]^{-(r+1)} \cdot [-(1-p)e^t]\\
&= p^r \cdot (-r) \cdot \left[1-\left(1-p\right)e^0\right]^{-(r+1)} \cdot [-(1-p)e^0]\\
&= p^r \cdot r \cdot p^{-(r+1)} \cdot (1-p)\\
&= r\frac{(1-p)}{p}
\end{align*}
\noindent
Variance of Negative Binomial is:

\begin{align*}
Var(X) &= \frac{d^2}{dt^2} \, \ln\left[M(0)\right]\\
&= \frac{d^2}{dt^2} \, \left[r\ln(p) - r \ln(1-(1-p)e^t)\right]\\
&= \frac{d}{dt} \, \left[ \frac{r(1-p)e^t}{1-(1-p)e^t}\right]\\
&= \frac{r(1-p)e^t[1-(1-p)e^t]-r(1-p)e^t[-(1-p)e^t]}{[1-(1-p)e^t]^2}\\
&= \frac{r(1-p)e^0[1-(1-p)e^0]+r(1-p)e^0[(1-p)e^0]}{[1-(1-p)e^0]^2}\\
&= r\frac{(1-p)}{p^2}
\end{align*}
\section*{3.1.27}
By definition we have:
\begin{align*}
    P(X \geq k+j \mid X \geq k) &= \frac{P(X \geq k+j \cap X \geq k)}{P(X \geq k)}\\
    &= \frac{P(X \geq k+j)}{P(X \geq k)}\\
    &= \frac{\sum_{i=k+j}^{\infty} (1-p)^{i}p}{\sum_{m=k}^{\infty} (1-p)^{m}p}\\
    &= \frac{(1-p)^{k+j}}{(1-p)^{k}}\\
    &= (1-p)^j\\
    &=\sum_{i=j}^{\infty} (1-p)^{i}p\\
    &= P(X \geq j)
\end{align*}
\section*{3.2.17}
Because \(X_1,X_2\) are independent and \(Y = X_1 + X_2\). Then we have 
\(M_Y(t) = M_{X_1}(t)M_{X_2}(t)\).
\begin{align*}
M_{X_2}(t) &= \frac{M_Y(t)}{M_{X_1}(t)}\\
&= \frac{exp[\mu (e^t-1)]}{exp[\mu_1 (e^t-1)]}\\
&= exp[(\mu - \mu_1)(e^t-1)]
\end{align*}

\noindent
By mgf of \(X_2\) and \(\mu > \mu_1\). We can identify that \(X_2 \sim Possion(\mu - \mu_1)\).
\section*{3.3.16}
\begin{align*}
P(Y<y) &= P(-2\ln(X)<y)\\
&=P(\ln(X)>-\frac{y}{2})\\
&=P(X>e^{-y/2})\\
&=\int_{e^{-y/2}}^{1} 1 \, dx\\
&=1-e^{-y/2}
\end{align*}
\begin{align*}
\text{Then when } X=0, Y= \infty \text{ and } X=1, Y=0
\end{align*}
Therefore, we have cdf of Y:
\begin{align*}
F_Y(y)=
\begin{cases}
1-e^{-y/2} \text{  if  } y \in (0, \infty)\\
0 \text{  elsewhere}
\end{cases}
\end{align*}
\begin{align*}
\frac{d}{dy}F_Y(y)&=\frac{d}{dy}1-e^{-y/2}\\
&=\frac{1}{2}e^{-y/2}
\end{align*}
The pdf of Y:
\begin{align*}
f_Y(y)=
\begin{cases}
\frac{1}{2}e^{-y/2} \text{  if  } y \in (0, \infty)\\
0 \text{  elsewhere}
\end{cases}
\end{align*}
\section*{3.3.24}
\subsection*{(a)}
Because \(X_1,X_2\) are independent.
\begin{align*}
M_Y(t)&=\mathbb{E}[e^{tY}]\\
&=\mathbb{E}[e^{t(2X_1+6X_2)}]\\
&=\mathbb{E}[e^{2tX_1}]\mathbb{E}[e^{6tX_2}]\\
&=M_{X_1}(2t)M_{X_2}(6t)\\
&=[1-3(2t)]^{-3} \cdot [1-1(6t)]^{-5}\\
&=(1-6t)^{-8}
\end{align*}
\subsection*{(b)}
By mgf of Y,we can know that \(Y \sim Gamma(\alpha = 8,\beta = 6)\).
\section*{3.3.25}
\subsection*{(a)}
\begin{align*}
P(X>x+y \mid X>x)&=\frac{P(X > x+y \cap X>x)}{P(X>x)}\\
&=\frac{P(X > x+y)}{P(X>x)}\\
&=\frac{1-F_X(x+y)}{1-F_X(x)}\\
&=\frac{e^{-\lambda(x+y)}}{e^{-\lambda x}}\\
&=e^{-\lambda y}\\
&=1-F_X(y)\\
&=P(X>y)
\end{align*}
\subsection*{(b)}
Let \(g(y) = 1 - F_Y(y) = P(Y > y)\).The memoryless property states:
\[
P(Y > x + y \mid Y > x) = P(Y > y)
\]
Using the definition of conditional probability:
\[
\frac{P(Y > x + y \cap Y > x)}{P(Y > x)} = 
\frac{P(Y > x + y)}{P(Y > x)} =
P(Y > y)
\]
Substituting $g(y)$:
\[
\frac{g(x + y)}{g(x)} = g(y) \implies g(x + y) = g(x)g(y)
\]
This is Cauchy's functional equation. Since $Y$ is a continuous random variable, 
$F_Y(y)$ is continuous, and thus $g(y)$ must be continuous. 
The only continuous solutions for $y > 0$ are of the form $g(y) = c^y$ 
for some constant $c$.\\
We are given $0 < F_Y(y) < 1$ for $y > 0$, 
which implies $0 < g(y) < 1$ for $y > 0$.This restricts the constant $c$ to 
$0 < c < 1$.\\
We can define a positive constant $\lambda = -\ln(c)$ (which is positive since $0 < c < 1$).
Then $c = e^{-\lambda}$.\\
So, $g(y) = c^y = (e^{-\lambda})^y = e^{-\lambda y}$.\\
Finally, we substitute back $g(y) = 1 - F_Y(y)$:
$$1 - F_Y(y) = e^{-\lambda y}$$
$$F_Y(y) = 1 - e^{-\lambda y}, \quad \text{for } y > 0$$
This is the cdf of an exponential distribution with parameter $\lambda$. 
We also verify $F_Y(0) = 1 - e^0 = 0$, which matches the given condition.
\section*{3.4.29}
Because $X_1$ and $X_2$ are independent with normal distributions $N(6, 1)$ and
$N(7, 1)$.\\
Therefore, we have:
\[
X_1-X_2 \sim N(6-7,1+1) \sim N(-1,2)
\]
Then,
\begin{align*}
P(X_1>X_2) &= P(X_1-X_2 >0)\\
&=1-\Phi(\frac{0-(-1)}{\sqrt{2}})\\
&=\Phi(-\frac{1}{\sqrt{2}})\\
&\approx 0.23975
\end{align*}
\section*{3.5.7}
Because X and Y have a bivariate normal distribution, we have $Y|X=x$:
\[
Y|X=x \sim N(\mu_2 - \rho (x - \mu_1), \sigma_2^2(1-\rho^2))
\sim N(10 - \rho (x - 5), 25(1-\rho^2))
\]
Then,
\begin{align*}
P(4<Y<16 \mid X=5)&=\Phi\left(\frac{16-\rho (5 - 5) -10}{\sqrt{25(1-\rho^2)}}\right)-
\Phi\left(\frac{4-\rho (5 - 5) -10}{\sqrt{25(1-\rho^2)}}\right)\\
&=\Phi\left(\frac{6}{5\sqrt{1-\rho^2}}\right)-
\Phi\left(\frac{-6}{5\sqrt{1-\rho^2}}\right)\\
&=2\Phi\left(\frac{6}{5\sqrt{1-\rho^2}}\right) -1
\end{align*}
And we know that
$$P(4<Y<16 \mid X=5) =  0.954 \approx 2\Phi(2)-1$$
Therefore
\[
\frac{6}{5\sqrt{1-\rho^2}}\approx 2 \Rightarrow \rho \approx 0.8
\]
\section*{3.5.12}
Because X and Y have a bivariate normal distribution with parameters $\mu_1 =\mu_2 =0$, 
$\sigma_1^2 = \sigma_2^2 = 1$,
and correlation coefficient $\rho$. 
\[
Z = aX+bY \sim N(a\mu_1+b\mu_2,a^2\sigma_1^2+b^2\sigma_2^2+2ab\rho\sigma_1\sigma_2)
\sim N(0,a^2+b^2+2ab\rho)
\]
\section*{3.6.11}
A random variable $F$ follows an F-distribution with $r_1$ and $r_2$ degrees of freedom, 
denoted $F \sim F(r_1, r_2)$, if it can be written as the ratio of two independent chi-square 
($\chi^2$) random variables, each divided by its respective degrees of freedom:
\[
F = \frac{U_1 / r_1}{U_2 / r_2}
\]
Where: $U_1 \sim \chi^2_{r_1}\, , \,  U_2 \sim \chi^2_{r_2}\, , \, U_1$ and $U_2$ are independent.
Then
\[
T^2 = \frac{W^2 / 1}{V / r}
\]
\begin{enumerate}
    \item The Numerator ($U_1 / r_1$):
    \begin{itemize}
        \item Let $U_1 = W^2$ and $r_1 = 1$.
        \item We are given $W \sim N(0, 1)$.
        \item By definition, the square of a standard normal random variable is a chi-square 
        random variable with 1 degree of freedom.
        \item Therefore, $U_1 = W^2 \sim \chi^2_1$.
        \item This matches the requirement for the numerator, with $r_1 = 1$.
    \end{itemize}
    \item The Denominator ($U_2 / r_2$):
    \begin{itemize}
        \item Let $U_2 = V$ and $r_2 = r$.
        \item We are given $V \sim \chi^2_r$.
        \item This matches the requirement for the denominator, with $r_2 = r$.
    \end{itemize}
    \item Independence:
    \begin{itemize}
        \item We need $U_1$ and $U_2$ to be independent.
        \item $U_1 = W^2$ and $U_2 = V$.
        \item We are given that $W$ and $V$ are independent. 
        Since $W^2$ is a function of $W$, it is also independent of $V$.
        \item Therefore, $U_1$ and $U_2$ are independent.
    \end{itemize}
\end{enumerate}
We have successfully shown that $T^2$ can be expressed as:
$$T^2 = \frac{W^2 / 1}{V / r} = \frac{U_1 / r_1}{U_2 / r_2}$$
Where: $U_1 = W^2 \sim \chi^2_1$ (so $r_1 = 1$) , $U_2 = V \sim \chi^2_r$ (so $r_2 = r$),
$U_1$ and $U_2$ are independent.
This is the exact definition of an F-distribution with $r_1 = 1$ and $r_2 = r$.
Thus, $T^2 \sim F(1, r)$.
\end{document}